\documentclass[10pt,a4paper]{amsart}
\usepackage[margin=19mm]{geometry}
\usepackage{amsmath,amssymb,mathtools}
\usepackage[scr=rsfs,cal=euler]{mathalfa}
\usepackage{xfrac}
\usepackage[unicode,hidelinks,bookmarks=false]{hyperref}
\newcommand{\ie}{i.e.\ }
\newcommand{\cf}{cf.\ }
\newcommand{\resp}{respectively\ }
\newcommand{\Subsection}{\subsection}
\providecommand{\nobreakdash}{\nobreak\hbox{-}}
\setlength{\emergencystretch}{2em}
\multlinegap0pt
\usepackage{mathtools}
\usepackage{amssymb}
\usepackage[shortlabels]{enumitem}
\newtheorem{lemm}{Lemma}
\newtheorem{theo}{Theorem}
\newtheorem{coro}{Corollary}
\newtheorem{prop}{Proposition}
\def\Vect{{\operatorname{Vect}}}
\def\band{{\mathfrak b}}
\def\cA{{\mathcal A}}
\def\cG{{\mathcal G}}
\def\cL{{\mathcal L}}
\def\cX{{\mathcal X}}
\def\ed{{\operatorname{ed}}}
\def\fC{{\mathfrak C}}
\def\ra{{\rightarrow}}
\def\rdim{{\operatorname{rdim}}}
\title{Essential dimension and faithful rank of finite $p$-gerbes}
\author{Tianzhi Yang}
\date{Source: arXiv:2609.01932v1 (CC BY 4.0)}
\begin{document}
\maketitle

\noindent\emph{Source and licence.} Essential dimension and faithful rank of finite $p$-gerbes. Version arXiv:2609.01932v1 (CC BY 4.0), distributed by arXiv under the Creative Commons Attribution 4.0 International licence, http://creativecommons.org/licenses/by/4.0/, as recorded in the arXiv licence metadata for this exact version and shown on the abstract page https://arxiv.org/abs/2609.01932v1. Full licence notice: \texttt{licenses/arxiv-2609.01932-CC-BY-4.0.txt}.\par\smallskip

\noindent\textit{Editorial scope.} Counted unit: the article's prop pro:generic-embedding-problem (source0.tex lines 2180--2210). The article's complete original proof of it is delivered with it (source0.tex lines 2212--2295, one \texttt{proof} environment of 84 lines). No mathematics is written, repaired or re-derived: the statement and the proof are the article's own lines, in the article's own notation, and no retained line is substituted or re-flowed.  Retained uncounted as the source-local context the proof needs: lines 1093--1099; lines 1483--1490; lines 1588--1594; lines 1679--1686; lines 2106--2109; lines 2129--2139.  Nothing was omitted from any retained span, so the package carries no omission record.  No retained line contains a citation, so the wrapper carries no bibliography and no bibliography entry.  No retained line contains a figure environment, an image-inclusion command or drawing code, so no image is delivered and the package declares no asset dependency.  Editorial formatting only, in addition: an AMS \texttt{amsart} preamble that restates the article's own declarations and macros used by the retained lines, namely the environments \texttt{lemm}, \texttt{theo}, \texttt{coro}, \texttt{prop} on separate counters, together with the standard packages those lines need.  The retained environments print in this document as Lemm 1, Theo 1, Coro 1, Proposition 1 in order of appearance, whereas the article prints the same items with the numbering of its own full document.  The complete native source of the article as distributed in the arXiv e-print for this exact version is bundled unchanged as \texttt{sources/209154/source0.tex}.
\begin{lemm}\label{lem:generic-central-g0-surjection}
For every $\chi\in\fC^*$, restriction along $\cX_\eta\to\cG$ induces a
surjection
\[
G_0^\chi(\cG)\twoheadrightarrow G_0^\chi(\cX_\eta).
\]
\end{lemm}

\begin{lemm}\label{lem:p-power-ranks-main}
Assume that $\mu_p\subset k$ and that $\band$ has $p$-primary
descent.  Every simple object of
$\Vect(\cG)$ has rank a power of $p$.  Consequently,
\begin{equation}\label{eq:block-gcd}
m_\chi(\cG)=d_\chi(\cG).
\end{equation}
\end{lemm}

\begin{lemm}\label{lem:ed-base-change}
For every field extension $K/k$ and every category fibered in groupoids
$\cA/k$,
\[
\ed_k(\cA;p)\ge\ed_K(\cA_K;p).
\]
\end{lemm}

\begin{theo}
\label{thm:gerby-km-special}
Assume that $\mu_p\subset k$ and that the band $\band$ has
$p$-primary descent.  Then
\begin{equation}\label{eq:gerby-km-special}
\ed_k(\cG)=\ed_k(\cG;p)=\rdim(\cG).
\end{equation}
\end{theo}

\begin{equation}\label{eq:p-group-extension}
1\ra N\ra P'
\xrightarrow{\pi}P\ra1
\end{equation}

\begin{coro}
\label{cor:embedding-problem}
For the extension \eqref{eq:p-group-extension},
\[
\sup_{\substack{K/k\\E\in H^1(K,P)}}
\ed_K\bigl(\operatorname{Lift}_{P'}(E);p\bigr)
=\rdim_p(P',N).
\]
Hence $\rdim_p(P',N)$ is the optimal uniform bound for the essential
$p$-dimension of the solution gerbes of this embedding problem.
\end{coro}

\begin{prop}
\label{pro:generic-embedding-problem}
Choose a $P$-representation $W$ with a nonempty $P$-stable open subset
$U\subset W$ on which $P$ acts freely.  Put
\[
K=k(U/P),
\qquad
E_{\mathrm{gen}}=U_K.
\]
Then
\begin{equation}\label{eq:generic-embedding-attainment}
\ed_K\bigl(\operatorname{Lift}_{P'}(E_{\mathrm{gen}});p\bigr)
=\rdim_p(P',N).
\end{equation}
In particular, $E_{\mathrm{gen}}$ realizes the supremum in
Corollary~\ref{cor:embedding-problem}.

Moreover, let $F=k^{(p)}$ be a $p$-closure of $k$, put
\[
K_F=F(U_F/P),
\]
and let $E_{\mathrm{gen},F}$ be the generic fiber of
$U_F\to U_F/P$.  Then
\begin{equation}\label{eq:generic-embedding-attainment-pclosure}
\ed_{K_F}\bigl(\operatorname{Lift}_{P'}(E_{\mathrm{gen},F})\bigr)
=
\ed_{K_F}\bigl(\operatorname{Lift}_{P'}(E_{\mathrm{gen},F});p\bigr)
=
\rdim_p(P',N).
\end{equation}
\end{prop}

\begin{proof}
We first prove the stronger assertion after passage to the $p$-closure.
Thus replace $k$ by $F=k^{(p)}$.  Then $k$ is $p$-special, hence
$\mu_p\subset k$.  After this replacement write again
\[
K=k(U/P),
\qquad
E_{\mathrm{gen}}=U_K,
\]
and set
\[
\cL=[E_{\mathrm{gen}}/P'],
\quad C=N\cap Z(P')[p].
\]
For $\chi\in C^*$, restriction induces a surjection
\begin{equation}\label{eq:generic-lifting-g0}
G_0^\chi(BP')\ra G_0^\chi(\cL).
\end{equation}
Indeed, it factors as
\[
G_0^\chi(BP')
\xrightarrow{\sim}G_0^\chi([W/P'])
\ra G_0^\chi([U/P'])
\ra G_0^\chi(\cL).
\]
The first arrow is homotopy invariance and the second is localization.
For the third, spread a coherent sheaf from the generic fiber over the
inverse image of some nonempty open subset of $U/P$, and apply localization,
exactly as in the proof of
Lemma~\ref{lem:generic-central-g0-surjection}.  Since restriction preserves rank,
\eqref{eq:generic-lifting-g0} gives
\begin{equation}\label{eq:generic-lifting-rank-ideals}
I_\chi(BP')=I_\chi(\cL).
\end{equation}

The extension $k(U)/K$ is $P$-Galois and neutralizes $\cL$.  Hence
the descent action on its geometric band $N$ contains the full
conjugation action of $P=P'/N$.  The kernel of a vector bundle on
$\cL$ is therefore represented by a subgroup of $N$ normal in
$P'$.  Every nontrivial such subgroup meets $C$.  Thus faithfulness on
$\cL$, and relative faithfulness for $BP'\to BP$, are both equivalent
to faithfulness on $C$.

In both character categories, every simple object has rank a power of $p$; hence the associated rank ideals are realized by least blocks, as in the proof of Lemma~\ref{lem:p-power-ranks-main}.  Using Equation \eqref{eq:generic-lifting-rank-ideals} together with the two minimum basis formulas, we obtain
\[
\rdim(\cL)=\rdim(BP'\to BP)=\rdim_p(P',N).
\]
Moreover, since the band of $\cL$ admits $p$-primary descent, Theorem~\ref{thm:gerby-km-special} yields
\[
\ed_K(\cL)=\ed_K(\cL;p)=\rdim(\cL)=\rdim_p(P',N).
\]
Thus \eqref{eq:generic-embedding-attainment-pclosure} follows.

We now return to the original base field $k$.  Put
\[
\cL_0=\operatorname{Lift}_{P'}(E_{\mathrm{gen}}).
\]
After base change to $F=k^{(p)}$, the generic torsor
$E_{\mathrm{gen}}$ becomes $E_{\mathrm{gen},F}$ over $K_F$, and hence
\[
(\cL_0)_{K_F}\simeq
\operatorname{Lift}_{P'}(E_{\mathrm{gen},F}).
\]
By Lemma~\ref{lem:ed-base-change} and the equality just proved,
\[
\ed_K(\cL_0;p)
\ge
\ed_{K_F}\bigl((\cL_0)_{K_F};p\bigr)
=
\rdim_p(P',N).
\]
On the other hand, Corollary~\ref{cor:embedding-problem} gives
\[
\ed_K(\cL_0;p)\le \rdim_p(P',N).
\]
Therefore
\[
\ed_K\bigl(\operatorname{Lift}_{P'}(E_{\mathrm{gen}});p\bigr)
=\rdim_p(P',N),
\]
which is \eqref{eq:generic-embedding-attainment}.  The assertion that
$E_{\mathrm{gen}}$ realizes the supremum now follows from
Corollary~\ref{cor:embedding-problem}.
\end{proof}

\end{document}
